\documentclass[10pt]{article}
% Math Packages
\usepackage{amsmath, mathtools}
\usepackage{amssymb}
\usepackage{amsthm}
\usepackage{amsfonts}
\usepackage{bbm}
\usepackage{breqn}
\usepackage[margin=1in]{geometry}
\usepackage{graphicx}
\usepackage{tikz}
\usetikzlibrary{arrows.meta}
\usetikzlibrary{calc}
\usepackage{forest}
\usepackage{tikz-qtree}
\graphicspath{ {./images/} }
\usepackage{hyperref}
\usepackage[capitalize]{cleveref}
\usepackage[shortlabels]{enumitem}
\usetikzlibrary{arrows,matrix,positioning}
\usepackage{multicol}

% for the pipe symbol
\usepackage[T1]{fontenc}

% Citing theorems by name. (source: https://tex.stackexchange.com/questions/109843/cleveref-and-named-theorems)
\makeatletter
\newcommand{\ncref}[1]{\cref{#1}\mynameref{#1}{\csname r@#1\endcsname}}

\def\mynameref#1#2{%
  \begingroup
  \edef\@mytxt{#2}%
  \edef\@mytst{\expandafter\@thirdoffive\@mytxt}%
  \ifx\@mytst\empty\else
  \space(\nameref{#1})\fi
  \endgroup
}
\makeatother

% Colorful Notes
\usepackage{color} \definecolor{Red}{rgb}{1,0,0} \definecolor{Blue}{rgb}{0,0,1}
\definecolor{Purple}{rgb}{.5,0,.5} \def\red{\color{Red}} \def\blue{\color{Blue}}
\def\gray{\color{gray}} \def\purple{\color{Purple}}
\newcommand{\rnote}[1]{{\red [#1]}} % \rnote{foo} gives '[foo]' in red
\newcommand{\pnote}[1]{{\purple [#1]}} % \pnote{foo} gives '[foo]' in purple
\newcommand{\bnote}[1]{{\blue #1}} % \bnote{foo} gives 'foo' in blue
\newcommand{\gnote}[1]{{\gray #1}} % \gnote{foo} gives 'foo' in gray
\newcommand{\Max}[1]{{\purple [#1]}} % \bnote{foo} then 'foo' is blue


% Claim numbering (the counter restarts after each proof environment)
\newcounter{claimcount}
\setcounter{claimcount}{0}
\newenvironment{claim}{\refstepcounter{claimcount}\par\addvspace{\medskipamount}\noindent\textbf{Claim \arabic{claimcount}:}}{}
\usepackage{etoolbox}
\AtBeginEnvironment{proof}{\setcounter{claimcount}{0}}
\newenvironment{claimproof}{\par\addvspace{\medskipamount}\noindent\textit{Proof of Claim  \arabic{claimcount}.}}{\hfill\ensuremath{\qedsymbol} \tiny{Claim}

  \medskip}
% Add claim support to cleverref
\crefname{claimcount}{Claim}{Claims}


% Math Environments
\newtheorem{theorem}{Theorem}
\newtheorem{assumption}[theorem]{Assumption}
\newtheorem{lemma}[theorem]{Lemma}
\newtheorem{proposition}[theorem]{Proposition}
\newtheorem{corollary}[theorem]{Corollary}
\newtheorem{question}[theorem]{Question}
\theoremstyle{definition}
\newtheorem{definition}[theorem]{Definition}
\newtheorem{remark}[theorem]{Remark}
\newtheorem{example}[theorem]{Example}
\newtheorem{notation}[theorem]{Notation}
\newtheorem{problem}[theorem]{Problem}

% Matrices and Column Vectors. 
\usepackage{stackengine}
\setstackgap{L}{1.0\normalbaselineskip}
\usepackage{tabstackengine}
\setstackEOL{;}% row separator
\setstackTAB{,}% column separator
\setstacktabbedgap{1ex}% inter-column gap
\setstackgap{L}{1.5\normalbaselineskip}% inter-row baselineskip
\let\nmatrix\bracketMatrixstack  %Usage: \nmatrix{1,2,3\4,5,6}
\newcommand\cv[1]{\setstackEOL{,}\bracketMatrixstack{#1}} %usage: \cv{1,2,3}

% Custom Math Coqmmands
\newcommand{\vt}{\vskip 5mm} % vertical space
\newcommand{\fl}{\noindent\textbf} % first line
\newcommand{\Fl}{\vt\noindent\textbf} % first line with space above
\newcommand{\norm}[1]{\left\lVert#1\right\rVert} % norm
\newcommand{\pnorm}[1]{\left\lVert#1\right\rVert_p} % p-norm
\newcommand{\qnorm}[1]{\left\lVert#1\right\rVert_q} % q-norm
\newcommand{\1}[1]{\textbf{1}_{\left[#1\right]}} % indicator function 
\def\limn{\lim_{n\to\infty}} % shortcut for lim as n-> infinity
\def\sumn{\sum_{n=1}^{\infty}} % shortcut for sum from n=1 to infinity
\def\sumkn{\sum_{k=1}^{n}} % shortcut for sum from k=1 to n
\def\sumin{\sum_{i=1}^{n}} % shortcut for sum from i=1 to n
\def\SAs{\sigma\text{-algebras}} % shortcut for $\sigma$-algebras
\def\SA{\sigma\text{-algebra}} % shortcut for $\sigma$-algebra
\def\Ft{\mathcal{F}_t} % time-indexed sigma-algebra (t)
\def\Fs{\mathcal{F}_s} % time-indexed sigma-algebra (s)
\def\F{\mathcal{F}} % sigma-algebra
\def\G{\mathcal{G}} % sigma-algebra
\def\R{\mathbb{R}} % Real numbers
\def\N{\mathbb{N}} % Natural numbers
\def\Z{\mathbb{Z}} % Integers
\def\E{\mathbb{E}} % Expectation
\def\P{\mathbb{P}} % Probability
\def\Q{\mathbb{Q}} % Q probability
\def\dist{\text{dist}} %Text 'dist' for things like 'dist(x,y)'
\newcommand{\indep}{\perp \!\!\! \perp}  %independence symbol
\def\Var{\mathrm{Var}} % Variance
\def\tr{\mathrm{tr}} % trace

% Brackets and Parentheses
\def\[{\left [}
    \def\]{\right ]}
% \def\({\left (}
%   \def\){\right )}



\usepackage{color}
\definecolor{Red}{rgb}{1,0,0}
\definecolor{Blue}{rgb}{0,0,1}
\definecolor{Purple}{rgb}{.5,0,.5}
\def\red{\color{Red}}
\def\blue{\color{Blue}}
\def\gray{\color{gray}}
\def\purple{\color{Purple}}
\definecolor{RoyalBlue}{cmyk}{1, 0.50, 0, 0}
\newcommand{\dempfcolor}[1]{{\color{RoyalBlue}#1}} 
\newcommand{\demph}[1]{\dempfcolor{{\sl #1}}}

% comment exactly one of the following line to show / hide the solutions
% \newcommand{\solution}[1]{{\purple #1}} % uncomment to show the solutions
\newcommand{\solution}[1]{} % uncomment to hide the solutions



\title{Lecture Notes for Math 372: \\Elementary Probability and Statistics}
\date{Last updated: \today}
% \author{mh}

\begin{document}
\begin{center}
  \section*{Math 372: Homework 12}
  \textit{Due Monday, December 1}
\end{center}

\textit{Instructions: Do not write your solutions on this sheet---you don't have enough space to show
your work. So use separate sheets of paper.}


\begin{problem}[Exercise 8.1.13]
  The calibration of a scale is to be checked by weighing a 10kg brick 25
  times. Suppose that the results of different weighings are independent of
  one another and that the weight on each trial is normally distributed with
  $\sigma=.2$ kg. Let $\mu$ denote the true average weight reading on the
  scale.
  \begin{enumerate}[(a)]
    \item What hypotheses should be tested?
    \item With the sample mean itself as the test statistic, what is the
    $p$-value when $\overline{x}=9.85$, and what would you conclude at
    significance level $0.01$?
    \item For a test with $\alpha=.01$, what is the probability that
    recalibration is judged unnecessary when in fact $\mu=10.1$? When
    $\mu=9.8$?
  \end{enumerate}
\end{problem}

\begin{problem}[Exercise 8.2.19]
  The melting point of each of 16 samples of vegetable oil was determined,
  resulting in $\overline{x}=94.32$. Assume that the distribution of the
  melting point is normal with $\sigma=1.20$
  \begin{enumerate}[(a)]
    \item Test $H_{0}:\mu=95$ versus $H_{a}:\mu\neq 95$ using a two-tailed
    level $.01$ test.
    \item If a level $0.01$ test is used, what is $\beta(94)$, the probability
    of a type-II error when $\mu=94$?
    \item What value of $n$ is necessary to ensure that $\beta(94)=.1$ when $\alpha=.01$?
  \end{enumerate}
\end{problem}

\begin{problem}[Logistic]
  The density function
  \begin{equation*}
    f(x) = \frac{e^{3-x}}{(1+e^{3-x})^{2}}
  \end{equation*}
  is an example of a \demph{logistic} distribution.
  \begin{enumerate}[(a)]
    \item Verify that $f$ is a pdf
    \item Find $\P\left[3 \leq X \leq 4\right]$
    \item Graph $f$. What does the mean appear to be? What about the median?
  \end{enumerate}
\end{problem}



\begin{problem}
  A balanced coin is flipped 400 times. Determine the number $x$ such that
  with probability approximately 0.80, the number of heads is between $200-x$
  and $200+x$.
\end{problem}


\begin{problem}
  A spinner from a board game randomly indicates a real number between 0 and
  10. The spinner is fair in the sense that it indicates a number in a given
  interval of with the same probability as it indicates a number in any other
  interval of the same length.
  \begin{enumerate}[(a)]
    \item Explain why the function
    \begin{equation*}
      f(x) = \left\{ \begin{array}{l@{\quad:\quad}l} 0.1 & 0 \leq x \leq 10\\
          0& \text{else}\end{array}\right.
    \end{equation*}
    is a probability density function for the spinner's values.
    \item What does your intuition tell you about the value of the mean? Check
    your guess by evaluating an integral.
  \end{enumerate}
\end{problem}



\begin{problem}[Central Limit Theorem]
  One hundred $6$-sided dice are rolled, and their numbers are added up. Use
  the central limit theorem to estimate the probability that the sum of the
  numbers on the dice is between $330$ and $380$.
\end{problem}

\begin{problem}
  % Adapted from Stewart's Calculus.
  The Hydrogen atom is composed of one proton in the nucleous and one
  electron, which moves about the nucleus. In the quantum theory of atomic
  structure, it is assumed that the electron does not move in a well-defined
  orbit. Instead, it occupies a state known as an \demph{orbital} which may be
  thought of as a ``cloud'' of negative charge surrounding the nucleous. At the
  state of lowest energy, called the \demph{ground state}, the shape of this
  cloud is assumed to be a sphere centered at the nucleus. This sphere is
  described in terms of the probability density function
  \begin{equation*}
    f(r) = \frac{4}{a_{0}^{3}}r^{2}e^{-2r/a_{0}},
  \end{equation*}
  where $a_{0}$ is a physical constant known as the \demph{Bohr radius}, and
  $a_{0} \approx 5.29 \times 10^{-11} \ \mathtt{meters}.$
  For each $r>0$, let $P(r)$ denote the probability that the electron will be
  found within the sphere of radius $r$ meters centered at the nucleus. Then 
  \begin{equation*}
    P(r) = \int_{0}^{r}f(s) ds.
  \end{equation*}
  \begin{enumerate}[(a)]
    \item Verify that $f$ is a probability density function.
    \item Find $\lim_{r\to\infty} f(r)$.
    \item For what value of $r$ does $f(r)$ achieve its maximum value?
    \item Graph the density function. (I recommend using Desmos.)
    \item Find the probability that an electron will be within the sphere of
    radius $4a_{0}$.
    \item At a moment in time, the distance of the electron from the nucleus
    is a random variable, which we will call $R$. A good ``best guess'' for
    this distance is the expected value. Calculate the expected distance of
    the electron from the nucleus in the ground state of the hydrogen atom.
    \item The variance $\Var(R)$ measures \demph{uncertainty} about the
    electron's position. Compute $\Var(R)$.
    \item The momentum $M$ of the electron is another random variable whose
    density function is related to that of $R$ (in a manner which is beyond
    the scope of this class). The \demph{Heisenberg uncertainty principle}
    says that there is a constant $c>0$ such that
    \begin{equation*}
      \Var(R)\Var(M)\geq c
    \end{equation*}
    In one or two sentences, interpret this inequality.
  \end{enumerate}
\end{problem}


% \begin{problem}
%   The time in hours $T$ between streptococcal infection and display of a sore throat is a
%   random variable whose probabililty density function is approximately
%   \begin{equation*}
%     f(t) 
%     = \left\{ 
%       \begin{array}{l@{\quad:\quad}l} \frac{1}{15676} t^{2} e^{-0.05 t},  &0 \leq t \leq 150\\ 0& \text{else} \end{array}\right.
%   \end{equation*}
%   \begin{enumerate}[(a)]
%     \item What is the probability that an infected patient will display
%     symptoms within the first 48 hours?
%     \item What is the probability that an infected patient will not display
%     symptoms until after 36 hours?
%   \end{enumerate}
% \end{problem}



\begin{problem}
  \label{problem:exercise-10.6}
  We are interested in testing whether a coin is balanced based on the number
  of heads $Y$ in 36 flips; that is, $H_{0}:p=.5$ vs $H_{a}:p\neq 5$. Suppose
  that we decide ahead of time to reject $H_{0}$ if $\left| Y-18 \right| \geq 4$.
  \begin{enumerate}[(a)]
    \item What is the value of $\alpha$?
    \item What is the value of $\beta$ if $p=.7$?
  \end{enumerate}
\end{problem}


\begin{problem}
  Please answer ``true'' or ``false'' to the following questions.
  \begin{enumerate}[(a)]
    \item The level of the test computed in Problem \ref{problem:exercise-10.6}(a) is
    the probability that $H_{0}$ is true.
    \item The value of $\beta$ in Problem \ref{problem:exercise-10.6}(b) is the
    probability that $H_{a}$ is true.
    \item In Problem \ref{problem:exercise-10.6}(b), $\beta$ was computed assuming
    the null hypothesis was false.
    \item If $\beta$ was computed when $p=.55$, the value would be larger than
    the value of $\beta$ obtained in Problem \ref{problem:exercise-10.6}(b).
    \item The probability that the test mistakenly rejects $H_{0}$ is $\beta$. 
  \end{enumerate}
\end{problem}



% \begin{problem}
%   The joint probability density function of $X$ and $Y$ is given by
%   \begin{equation*}
%     f(x,y)
%     = \left\{ \begin{array}{l@{\quad:\quad}l} \frac{27 y^{2} e^{-3x}}{x^{2}} &0<y<x<\infty \\
%         0 & \text{otherwise} \end{array}\right.
%   \end{equation*}
%   \begin{enumerate}[(a)]
%     \item Find the marginal density for $X$.
%     \item Are $X$ and $Y$ indpendent? Why or why not?
%     \item Find $\P\left[X>2Y \right]$
%   \end{enumerate}
% \end{problem}

\begin{problem}%[continuous r.v.]
  Let $T$ be the waiting time, in hours, between successive purchases at a
  furniture store. Suppose that $T$ has pdf
  \begin{equation*}
    f_{T}(t) = C e^{-\frac{t}{4}}\mathbf{1}_{\left[ t \geq 0\right] } = \left\{ \begin{array}{l@{\quad:\quad}l} Ce^{-\frac{t}{4}} & t \geq 0\\ 0& t<0\end{array}\right.
  \end{equation*}
  where $C>0$ is some number that you'll have to determine.
  \begin{enumerate}[(a)]
    \item Find a value of $C$ so that $f_{T}(t)$ is a valid pdf.
    \item What is the probability that the time between purchases is greater
    than $6$ minutes? [Note: $6\text{ minutes }= 1/10\text{ of an hour}$].
    \item Compute $\E\left[T \right]$.
  \end{enumerate}
\end{problem}

% \section{Transplanted}

% \begin{problem}[Method of Moments]
%   Let $X$ be the proportion of allotted time that a randomly-selected student
%   spends working on an exam. Suppose that $X$ is a random variable with pdf
%   \begin{equation*}
%     f(x) = \left\{ \begin{array}{l@{\quad:\quad}l} (\theta+1)x^{\theta}
%         &0 \leq x \leq 1 \\ 0& \text{otherwise} \end{array}\right.
%   \end{equation*}
%   where $\theta\in \R$ is an unknown parameter. A random sample of 10 students
%   yields data
%   \begin{gather*}
%     x_{1} = .92,\ x_{2} = .79,\ x_{3} = .90,\ x_{4}=.65,\ x_{5}=.86\ \\
%     x_{6}=.47,\ x_{7}=.73,\ x_{8}=.97,\ x_{9}=.94,\ x_{10}=.77
%   \end{gather*}
%   Use the method of moments to obtain an estimator of $\theta$ and then
%   compute the estimate for this data.
% \end{problem}


% % \begin{problem}
% %   Let $S_{200}$ be the number of heads that turn up in 200 tosses of a fair
% %   coin. Estimate
% %   \begin{equation*}
% %     \P\left[S_{200}=100 \right], \P\left[S_{200}=90 \right], \text{ and } \P\left[S_{200}=80 \right].
% %   \end{equation*}
% % \end{problem}



% \begin{problem}
%   Find the probability that among $10,000$ random digits the digit 3 appears
%   not more than 931 times.



% \begin{problem}
%   A random walker starts at 0 on the x-axis and at each time unit moves 1 step
%   to the right or 1 step to the left with probability 1/2. Estimate the
%   probability that, after 100 steps, the walker is more than 10 steps from the
%   starting position.
% \end{problem}




% \newpage
% \section*{Solution to Problem 6}

% Let $J$ be June's score, and let $A$ be April's score. We want to compare
% \begin{equation*}
%   \P\left[J \geq 30\right] \quad \text{and} \quad  \P\left[A \geq 30\right]
% \end{equation*}

% To do this, we will use the Central Limit Theorem.

% \Fl{Calculation for June:} First observe that we can write
% \begin{equation*}
%   J = X_{1}+\ldots+X_{48}
% \end{equation*}
% where $X_{1},X_{2},\ldots,X_{48}$ are IID random variables with
% \begin{equation*}
%   X_{i}= \left\{ \begin{array}{l@{\quad:\quad}l} 1 & \text{with probability
%     }3/4\\0 &\text{with probability }1/4 \end{array}\right.
% \end{equation*}
% In this case, we have
% \begin{equation*}
%   \mu = \E\left[X_{i} \right] = \frac{3}{4}
% \end{equation*}
% and
% \begin{equation*}
%   \sigma^{2} = \Var(X_{i}) = \frac{3}{4} - \left( \frac{3}{4} \right)^{2} = \frac{3}{4}\frac{1}{4} = \frac{3}{16},
% \end{equation*}
% so that
% \begin{equation*}
%   \sigma = \frac{\sqrt{3}}{4}
% \end{equation*}
% Therefore by the Central Limit Theorem, we have
% \begin{align*}
%   \frac{J}{48} 
%   &\approx \mu + \frac{\sigma}{\sqrt{48}}Z \\
%   &= \frac{3}{4}+ \frac{\sqrt{3}}{4 \sqrt{48}}Z,
% \end{align*}
% where $Z$ is a standard normal random variable. Therefore,
% \begin{align*}
%   \P\left[J \geq 30\right] 
%   &=\P\left[\frac{J}{48}\geq \frac{30}{48}\right]\\
%   &\approx\P\left[  \frac{3}{4}+ \frac{\sqrt{3}}{4 \sqrt{48}}Z \geq  \frac{30}{48}\right]\\
%   &= \P\left[ \left( \frac{\sqrt{3}}{4 \sqrt{48}} \right)Z \geq -\frac{1}{8}\right]\\
%   &= \P\left[Z \geq -2\right]\\
%   &\approx 0.98
% \end{align*}
% Therefore June is about 98\% likely to pass the exam.

% \Fl{Calculation for April} We can write
% \begin{equation*}
%   A = Y_{1}+\ldots+Y_{48}
% \end{equation*}
% where $Y_{1},\ldots,Y_{48}$ are IID random variables with
% \begin{equation*}
%   X_{i}= \left\{ \begin{array}{l@{\quad:\quad}l} 1 & \text{with probability
%     }1/2\\0 &\text{with probability }1/2 \end{array}\right.
% \end{equation*}
% By similar calculations, we have $\mu = \E\left[Y_{i} \right] = \frac{1}{2}$  and
% $\sigma^{2} = \Var(Y_{i}) = \frac{1}{4}$, so $\sigma = \frac{1}{2}$.
% Therefore by the Central Limit Theorem,
% \begin{align*}
%   \frac{A}{48} 
%   &\approx \mu + \frac{\sigma}{\sqrt{48}}Z\\
%   &= \frac{1}{2} + \frac{1}{2 \sqrt{48}}Z,
% \end{align*}
% where $Z$ is a standard normal random variable. Using this approximation, we have:
% \begin{align*}
%   \P\left[A \geq 30\right]
%   &=\P\left[\frac{A}{48}\geq \frac{30}{48}\right]\\
%   &\approx \P\left[\frac{1}{2}+\frac{1}{2 \sqrt{48}}Z \geq \frac{30}{48}\right]\\
%   &= \P\left[Z \geq \sqrt{3}\right]\\
%   &= 0.04
% \end{align*}
% \Fl{Conclusion:} The probability that April passes the exam is $0.04$. The
% probability that June will pass the exam is $0.98$.


% \newpage
% \section*{Solution to Problem 7}
% Let
% \begin{equation*}
%   X_{i} = \left\{ \begin{array}{l@{\quad:\quad}l} 1 & \text{if the i-th
%       random digit is 3} \\0 & \text{if the i-th digit is not 3} \end{array}\right.
% \end{equation*}
% Since there are 10 digits $(0,1,2,\ldots,9)$, we have
% \begin{equation*}
%   \P\left[X_{i} =1\right] = \frac{1}{10}.
% \end{equation*}
% Then a straightforward computation shows that
% \begin{align*}
%   \mu &= \E\left[X_{i} \right] = \frac{1}{10}\\
%   \sigma^{2}&= \frac{1}{10} \left( 1-\frac{1}{10} \right)= \frac{9}{100}\\
%   \sigma&= \frac{3}{10}
% \end{align*}
% Let $S_{n}=X_{1}+\ldots+X_{n}$ be the number of 3's in the first $n$ random
% digits. We want to compute $\P\left[S_{10,000}\leq 931\right]$. By the Central Limit Theorem,
% \begin{align*}
%   \frac{S_{n}}{n} 
%   &\approx \mu + \frac{\sigma}{\sqrt{n}}Z.
% \end{align*}
% where $Z$ is a standard normal random variable. Here, we have $n=10,000$ so
% this becomes
% \begin{equation*}
%   \frac{S_{10000}}{10000} \approx \frac{1}{10}+ \frac{3}{1000}Z
% \end{equation*}
% Thus
% \begin{align*}
%   \P\left[S_{10000}\leq 931\right]
%   &= \P\left[\frac{S_{10000}}{10000} \leq  \frac{931}{10000}\right]\\
%   &\approx \P\left[  \frac{1}{10}+ \frac{3}{1000}Z \leq \frac{931}{10000}\right]\\
%   &= \P\left[Z \leq -2.3\right]\\
%   &\approx 0.01
% \end{align*}
% In other words, the probability of having 931 or fewer 3's out of 10,000
% random digits is about 1\%.

% \section{New problems}



\begin{problem}%[Extra credit]
  An urn contains one red ball and one blue ball. A ball is drawn at random.
  If it is red, the game is over. If it is blue, it is replaced in the urn
  along with an extra blue ball. This process continues until the red ball is
  drawn, however long it takes. Let $Y$ be the number of draws required.
  \begin{enumerate}[(a)]
    \item Argue that $\P\left[Y=k \right] = \frac{1}{k(k+1)}$.
    \item Show that this defines a valid pmf. (Hint: first show that
    $\frac{1}{k(k+1)}= \frac{1}{k}-\frac{1}{k+1}$, and then consider
    $\sum_{k=1}^{n}\frac{1}{k(k+1)}$).
    \item Show that $\E\left[Y \right]$ does not exist.
  \end{enumerate}
\end{problem}


% \begin{problem}[Normal Distribution]
%   If $Z\sim \mathcal{N}(0,1)$, then
%   \begin{equation}\label{eq:1}
%     \P\left[|Z| \leq u \right] = 2 \P\left[Z \leq u\right]-1
%   \end{equation}
%   for any $u \geq 0$. Draw a picture to illustrate the geometric meaning of \cref{eq:1}.
% \end{problem}


\end{document}